% Research draft for R75; deliberately outside the sealed R74 source.
% Depends on the retained permutation-core normal form and positivity argument.
\section{Orbit multiplicity, positive synthesis, and the convexity boundary}
\label{sec:r75-orbit-multiplicity}

The quotient reduction of Section~\ref{sec:r74-quotient-memory} has a precise
boundary. A single transitive orbit admits convexification without enlarging
the orbit exactly when its stabilizer is normal. For general finite groups,
irreducible multiplicities measure the obstruction and determine the exact
number of prescribed orbit copies at sufficiently small positive contrast.
All representations below are finite dimensional. We use classical complete
reducibility, matrix-coefficient orthogonality, and Frobenius reciprocity
\cite{Serre77}; their representation-theoretic content is not claimed as new.

Let $G$ be a finite group, not necessarily the faithful observable group of
the particular table under discussion. For a nonempty finite right $G$-set
$X$, let $\mathcal R_X(C)$ denote the probability tables
\[
 b(g)=\lambda P_gD,\qquad
 \lambda\in\Delta(X),\quad D(x,\cdot)\in\Delta(C),
\]
where $P_gP_h=P_{gh}$. Thus the same preparation and decoder are used for
every group element. Taking an observable quotient does not change the
table; keeping $G$ explicit lets us compare different allowed actions.
For this ambient-group formulation, $\kappa_G(b)$ denotes the minimum
$|X|$ over all such factorizations; an action need not be faithful unless
$G$ is the observable group.
Write $\widehat G$ for the complex irreducible representations, choose a
unitary representative $\rho$ of degree $d_\rho$, and define
\begin{equation}\label{eq:r75-fourier-block}
 \widehat b_c(\rho)=\frac1{|G|}\sum_{g\in G}b(g,c)\rho(g^{-1}),
 \qquad
 A_\rho(b)=\begin{pmatrix}
 \widehat b_{c_1}(\rho)\\ \vdots\\ \widehat b_{c_{|C|}}(\rho)
 \end{pmatrix},
 \qquad r_\rho(b)=\operatorname{rank}_{\mathbb C}A_\rho(b).
\end{equation}
Changing unitary bases does not change these ranks.

\begin{lemma}[Fourier multiplicity and real descent]
\label{lem:r75-fourier-multiplicity}
Let $V_b$ be the real span of the translates $b^h(g)=b(hg)$. The
multiplicity of the right irreducible $\rho$ in $V_b\otimes_{\mathbb R}
\mathbb C$ is $r_\rho(b)$. In particular $r_{\boldsymbol 1}(b)=1$ and
$r_\rho(b)\le d_\rho$. If $U$ and $V$ are real $G$-modules, then $U$
embeds equivariantly in $V$ if and only if every complex irreducible
multiplicity in $U\otimes\mathbb C$ is at most the corresponding one in
$V\otimes\mathbb C$.
\end{lemma}
\begin{proof}
The Fourier transform sends $b^h$ to the blocks
$A_\rho(b)\rho(h)$. Matrix-coefficient orthogonality identifies
$\mathbb C[G]$ with $\bigoplus_{\rho\in\widehat G}M_{d_\rho}(\mathbb C)$:
the group algebra can independently prescribe an arbitrary matrix on each
irreducible block. Consequently the $\rho$-part of the cyclic module
generated by $b$ identifies with
\[
 \{A_\rho(b)M:M\in M_{d_\rho}(\mathbb C)\}.
\]
If $A_\rho$ has rank $r$, a choice of a basis of its column space
identifies this right module with $r$ independent rows of length
$d_\rho$. It therefore consists of $r$ copies of $\rho$. For the trivial
representation the stacked block is the nonzero probability vector
$\pi=|G|^{-1}\sum_gb(g)$, so its rank is one.

Necessity of the second assertion follows by complexification. For
sufficiency, complete reducibility over $\mathbb C$ and the multiplicity
inequalities give a complex equivariant injection
$F:U\otimes\mathbb C\to V\otimes\mathbb C$. In real bases write
$F=A+iB$. Both $A$ and $B$ are real intertwiners. A maximal nonzero minor
of $F$ is the value at $z=i$ of a real polynomial minor of $A+zB$.
That polynomial is not identically zero, so some real $z$ makes this
minor nonzero. The real map $A+zB$ is the required injection. One may
choose $z$ among $0,1,\ldots,\dim U$, since the nonzero polynomial has
degree at most $\dim U$. This proves the assertion without excluding
complex or quaternionic representation types.
\end{proof}

For $H\le G$ put
\[
 m_\rho(H)=\dim_{\mathbb C}\rho^H
 =\operatorname{rank}\left(\frac1{|H|}\sum_{h\in H}\rho(h)\right).
\]
By Frobenius reciprocity these are the multiplicities of
$\mathbb C[H\backslash G]$. Let $q_H(b)$ be the least positive integer
$q$ for which $b$ belongs to $\mathcal R_{q(H\backslash G)}(C)$, where
$q(H\backslash G)$ is the disjoint union of $q$ copies; set $q_H(b)=\infty$
if no such factorization exists. This is an action-restricted resource.
It is not $\kappa_G(b)/[G:H]$ in general, since an unrestricted optimal
core may use other orbit types.

\begin{theorem}[Exact orbit multiplicity at small contrast]
\label{thm:r75-orbit-multiplicity}
Let $b:G\to\Delta(C)$, let
$\pi=|G|^{-1}\sum_gb(g)$, and set $b_t=\pi+t(b-\pi)$ for $0\le t\le1$.
For $H\le G$ define
\begin{equation}\label{eq:r75-copy-number}
 Q_H(b)=\max_{\rho\in\widehat G}
 \left\lceil\frac{r_\rho(b)}{m_\rho(H)}\right\rceil,
\end{equation}
with a positive numerator and zero denominator interpreted as $\infty$
and $0/0$ as zero. For every $t>0$,
$q_H(b_t)\ge Q_H(b)$. If $Q_H(b)<\infty$, an explicitly positive $t_0$
exists such that
\[
 q_H(b_t)=Q_H(b)\qquad(0<t\le t_0).
\]
If $Q_H(b)=\infty$, no positive contrast is realizable using any finite
number of copies of this orbit. At zero contrast one copy suffices.

More generally, choose one representative of each subgroup conjugacy
class. The unrestricted permutation-module bound has the exact finite
integer formulation
\begin{equation}\label{eq:r75-multiplicity-cover}
 \nu_G(V_b)=
 \min\left\{
  \sum_H q_H[G:H]:\ q_H\in\mathbb Z_{\ge0},\quad
  \sum_Hq_Hm_\rho(H)\ge r_\rho(b)\quad(\rho\in\widehat G)
 \right\}.
\end{equation}
Consequently this integer minimum equals $\kappa_G(b_t)$ for every
sufficiently small positive contrast, with the pure-start charge given
by Theorem~\ref{thm:r73-memory-formula} when $G$ is the observable group.
\end{theorem}
\begin{proof}
A factorization through a finite $G$-set $Y$ makes $V_{b_t}$ a quotient
of a submodule of $\mathbb R^Y$. Averaging an inner product over the
finite group makes every invariant subspace have an invariant orthogonal
complement; hence this subquotient embeds into $\mathbb R^Y$. The
contrast change is the identity on the invariant line generated by
$\pi$ and multiplication by $t$ on the zero-average part, so for $t>0$
its module is isomorphic to $V_b$. On $q(H\backslash G)$ the available
complex multiplicities are $qm_\rho(H)$, proving the lower bound and
the assertion when an active type is absent.

For the matching construction, let $Y$ be any nonempty finite $G$-set
whose complex multiplicities dominate $r_\rho(b)$. The lemma gives a
real equivariant injection of $V_b$ into $\mathbb R^Y$. Here are the
positivity details, to make clear that no complex factorization is
being substituted for a stochastic one. Delete the output columns
with $\pi_c=0$; nonnegativity makes them identically zero, and they can
be restored afterwards. Decompose
\[
 V_b=\mathbb R\overline b\oplus W,\qquad
 \overline b(g)=\pi,
\]
where $W$ is the kernel of group averaging. Every element of $W$ has
zero total output mass. Put $k=|Y|$, let $p$ be its uniform probability
row, and choose an equivariant injection $j:W\to\mathbb R^Y$.
The image of $j$ is orthogonal to all invariant vectors. Let
$F:\mathbb R^Y\to W$ be orthogonal projection onto $j(W)$ followed by
$j^{-1}$. Then $Fj=I_W$ and $F(p)=0$. Set
\[
 a=j(b-\overline b),\qquad T(y,c)=F(e_y)(e,c).
\]
Equivariance gives $aP_gT=b(g)-\pi$, while $pT=0$ and each row of $T$
sums to zero. If $b$ is nonconstant, choose
\[
 \epsilon=\frac1{2k\|a\|_\infty},\qquad
 B=\max_{y,c}\frac{|T(y,c)|}{\pi_c},\qquad
 t_0=\min\{1,\epsilon/B\}.
\]
The probability row $\lambda=p+\epsilon a$ and the row-stochastic
decoder $D_t(y,c)=\pi_c+(t/\epsilon)T(y,c)$ then satisfy
$\lambda P_gD_t=b_t(g)$ for $0\le t\le t_0$. For a constant table the
constant decoder works at every contrast. This is precisely the
retained positive construction, applied to the specified action
rather than only to a globally minimum one. Taking
$Y=Q_H(b)(H\backslash G)$ proves equality in the copy formula.

Every finite $G$-set is a disjoint union of coset orbits, and each such
orbit contributes the multiplicities $m_\rho(H)$. The real descent
lemma therefore makes the displayed integer constraints equivalent to
real module containment. Conversely every feasible integer vector
defines that disjoint-union action. This proves
\eqref{eq:r75-multiplicity-cover}; Theorem~\ref{thm:r74-low-contrast}
or the preceding construction gives its positive realization. The
regular orbit makes the minimum finite, since
$m_\rho(\{e\})=d_\rho\ge r_\rho(b)$.
\end{proof}

\begin{theorem}[Normality is the exact single-orbit convexity boundary]
\label{thm:r75-normality-convexity}
For $H\le G$ and the transitive right action $X=H\backslash G$, the
following conditions are equivalent:
\begin{enumerate}
\item $H$ is normal in $G$;
\item $\mathcal R_X(\{0,1\})$ is convex;
\item $\mathcal R_X(C)$ is convex for every finite outcome set $C$.
\end{enumerate}
If $H$ is normal, this response set is exactly the polytope of
probability tables constant on the left cosets of $H$. If $H$ is not
normal, nonconvexity occurs in every neighborhood of the interior
constant binary table $(1/d,1-1/d)$, where $d=[G:H]$.

An explicit witness is supplied by the return tables
\[
 f_x(g)=\mathbf1_{\{xg=x\}},\qquad x\in X,
 \qquad
 F_X(g)=\frac1d\sum_{x\in X}f_x(g)=\frac{\chi_X(g)}d.
\]
Every $(f_x,1-f_x)$ lies in $\mathcal R_X(\{0,1\})$, whereas
$(F_X,1-F_X)$ does not when $H$ is not normal. The same conclusion
holds after replacing every scalar $f$ by $d^{-1}+t(f-d^{-1})$, for
every $0<t\le1$.
More quantitatively, for any irreducible $\rho$ with
$0<m_\rho(H)<d_\rho$, put
$F_{X,t}=d^{-1}+t(F_X-d^{-1})$. Then
\begin{equation}\label{eq:r75-nonconvexity-margin}
 \inf_{(a,1-a)\in\mathcal R_X(\{0,1\})}
       \max_{g\in G}|F_{X,t}(g)-a(g)|
 \ \ge\ \frac{t\,m_\rho(H)}{d\,d_\rho}>0.
\end{equation}
This is a certified separation margin; optimality of its constant
is not asserted.
\end{theorem}
\begin{proof}
If $H$ is normal, for
$b(g,c)=\sum_{Hx}\lambda(Hx)D(Hxg,c)$ one has
$b(hg,c)=b(g,c)$ for every $h\in H$: normality gives $Hxh=Hx$.
Conversely every probability table with this left invariance is
realized by the point preparation $H$ and decoder
$D(Hg,c)=b(g,c)$. This proves the stated polytope identity and
(i)$\Rightarrow$(iii)$\Rightarrow$(ii).

Suppose $H$ is not normal. The group-algebra element
$e_H=|H|^{-1}\sum_{h\in H}h$ is an orthogonal projection in every
unitary irreducible representation, of rank $m_\rho(H)$. If every one
of these ranks were either zero or $d_\rho$, each matrix $\rho(e_H)$
would be zero or the identity. The faithful decomposition of the
complex group algebra into its irreducible matrix blocks would then
make $e_H$ central. Its coefficients are constant under conjugation
exactly when $H$ is normal. Hence some irreducible $\rho$ has
\[
 0<m_\rho(H)<d_\rho.
\]
This irreducible is nontrivial.

The return table at $x$ is realized by initial point $x$ and the
indicator decoder of $x$. Since
$\chi_X=\sum_\sigma m_\sigma(H)\chi_\sigma$, orthogonality gives
\[
 \widehat F_X(\rho)=\frac{m_\rho(H)}{d\,d_\rho}I_{d_\rho}.
\]
Its rank is $d_\rho>m_\rho(H)$. Lemma~\ref{lem:r75-fourier-multiplicity}
and the module lower bound therefore exclude a single-orbit
factorization. This average of members of the response set proves
nonconvexity. For positive contrast, the nontrivial Fourier block is
multiplied by $t$ and keeps its rank. Each contrasted return table is
still realized on its original orbit by replacing its decoder with
$d^{-1}+t(D-d^{-1})$. Letting $t$ decrease to zero proves the
neighborhood assertion. Thus (ii) implies (i).

For the quantitative assertion let $A=\widehat a(\rho)$ for any
single-orbit response. It has rank at most $m_\rho(H)<d_\rho$,
so there is a unit vector $v$ with $Av=0$. The Fourier block of
$F_{X,t}$ is $\alpha I$, where
$\alpha=t\,m_\rho(H)/(d\,d_\rho)$. Therefore
$\|\widehat F_{X,t}(\rho)-A\|_{\mathrm{op}}\ge\alpha$.
Because every $\rho(g)$ is unitary, the definition of the normalized
Fourier transform also gives
\[
 \|\widehat F_{X,t}(\rho)-\widehat a(\rho)\|_{\mathrm{op}}
 \le\frac1{|G|}\sum_{g\in G}|F_{X,t}(g)-a(g)|
 \le\max_{g\in G}|F_{X,t}(g)-a(g)|.
\]
Combining these inequalities proves the margin for every allowed
probability preparation and bounded decoder.
\end{proof}

For this witness the exact small-contrast number of copies is
\begin{equation}\label{eq:r75-character-copies}
 Q_H(F_X,1-F_X)
 =\max_{\rho:m_\rho(H)>0}
   \left\lceil\frac{d_\rho}{m_\rho(H)}\right\rceil.
\end{equation}
Indeed every active Fourier block of the character is full rank,
including the trivial block. Thus nonnormality entails a strict copy
cost at arbitrarily small positive contrast, with a matching positive
realization. This obstruction is not caused by a decoder touching the
boundary of its simplex.

\begin{lemma}[Full-contrast return-law rigidity]
\label{lem:r75-return-rigidity}
Let $X$ be a transitive $d$-point $G$-set. If a scalar response is a
convex combination of responses on copies of $X$, has value one at
$e$, and has group average $1/d$, every positive-weight component
uses a point preparation and its point indicator decoder. In
particular every component is a return law $f_x$.
\end{lemma}
\begin{proof}
For a component write $\lambda\in\Delta(X)$ and $D\in[0,1]^X$.
The total value one at $e$ forces $\lambda D=1$ for every component
of positive weight, hence $D(x)=1$ on the nonempty support of
$\lambda$. Transitivity gives its group average
$d^{-1}\sum_xD(x)\ge1/d$. Equality of the weighted mean with $1/d$
forces equality for each component. Consequently $\sum_xD(x)=1$,
so the support of $\lambda$ consists of exactly one point and $D$
is its indicator.
\end{proof}

These statements separate three matters: the classical multiplicities
of a permutation representation, the positivity of its preparation
and decoder, and the cost of an externally prescribed pure start.
The integer formula is exact at small positive contrast; it does not
replace the general full-contrast bilinear synthesis problem by a
linear program. The convexity theorem pinpoints why the retained
normal-stabilizer absorption cannot simply be extended to arbitrary
subgroups.

% Joint root/child research derivation for R75. Not a sealed source edit.
\subsection{An unbounded nonabelian family with a strict positivity cost}
\label{sec:r75-symmetric-return}

We apply the multiplicity theory to all symmetric groups $S_n$ with
$n\ge5$. The natural permutation character $\chi_n(g)$ counts the
fixed points of $g$ on $[n]$. Consider the binary response whose first
coordinate is
\begin{equation}\label{eq:r75-symmetric-return}
 \beta_t(g)=\frac{1-t}{n}+\frac{t}{n}\chi_n(g),
 \qquad b_t=(\beta_t,1-\beta_t),\qquad0\le t\le1.
\end{equation}
Its group average is $1/n$, and at the identity it equals
$[1+(n-1)t]/n$. This supplies a family with growing group, observable
rank, and exact memory, rather than a fixed-dimensional illustration.

The positivity threshold is controlled by the following matrix
extremum. No symmetry or positive-semidefiniteness hypothesis is
imposed on its coefficient matrix.

\begin{lemma}[A sharp singular-coefficient bound and its equality case]
\label{lem:r75-singular-coefficient}
Let $n\ge5$, $P=I-\mathbf1\mathbf1^{\mathsf T}/n$, and $t>0$.
If an entrywise nonnegative real $n\times n$ matrix $A$ satisfies
\[
 \operatorname{rank}A<n,\qquad
 \sum_{i,j}A_{ij}=1,\qquad PAP=\frac tnP,
\]
then
\begin{equation}\label{eq:r75-sharp-threshold}
 t\le\tau_n:=\frac{\sqrt n-1}{\sqrt n+1}.
\end{equation}
Equality holds if and only if, after the same permutation of rows
and columns,
\begin{equation}\label{eq:r75-extremal-coefficient}
 A=A_*:=\frac1{n-1}\sum_{i=1}^{n-1}
                  \ell_i^{\mathsf T}\ell_i,
 \qquad
 \ell_i=\left(1-\frac1{\sqrt n}\right)e_i+
                         \frac1{\sqrt n}e_n.
\end{equation}
Here $e_i$ and $\ell_i$ are row vectors. In particular symmetry of
the extremal matrix is a conclusion.
\end{lemma}
\begin{proof}
Put $\alpha=t/n$. The centered identity implies
$A=\alpha I+x\mathbf1^{\mathsf T}+\mathbf1y^{\mathsf T}$:
indeed $P(A-\alpha I)P=0$ decomposes the difference into a matrix
constant on each row and one constant on each column. Let $v$ be a
nonzero left null vector of $A$. If $v^{\mathsf T}\mathbf1=0$,
then $v^{\mathsf T}AP=v^{\mathsf T}PAP=\alpha v^{\mathsf T}$,
a contradiction. Normalize $\sum_i v_i=1$. Substitution of
$v^{\mathsf T}A=0$ in the preceding decomposition gives a vector
$c$ such that
\begin{equation}\label{eq:r75-null-vector-form}
 A=\alpha\{I+c\mathbf1^{\mathsf T}-\mathbf1v^{\mathsf T}\},
 \qquad v\cdot c=0,
 \qquad 1=t\sum_i c_i.
\end{equation}
We bound $\sum_i c_i$. Relabel so that $v_1=M$ is maximal and
$v_2=L$ is second largest. Since $\sum_i v_i=1$, $M>0$.
Entrywise nonnegativity gives
\[
 c_1\ge\max\{L,M-1\},\qquad c_i\ge M\quad(i>1).
\]
If $v$ is nonnegative, then $M\le1$ and $L\ge0$.
Thus $v\cdot c=0$ forces $v_i=0$ for every $i>1$,
$v_1=1$, and $c_1=0$. Consequently $\sum_i c_i\ge n-1$.

Otherwise put $m=-\min_i v_i>0$ and
$b_0=\max\{L,M-1\}$. Write
$c=(b_0,M,\ldots,M)+u$ with $u\ge0$. The dot product of $v$
with the displayed base vector is $M(b_0+1-M)\ge0$.
Since $v\cdot u\ge-m\sum_i u_i$, the condition $v\cdot c=0$
therefore yields
\begin{equation}\label{eq:r75-null-sum-bound}
 \sum_i c_i\ge b_0+(n-1)M+
                       \frac{M(b_0+1-M)}m.
\end{equation}
The normalization of $v$ also implies
\[
 M\ge M_-:=\frac{1+m}{n-1},\qquad
 L\ge\frac{1-M+m}{n-2}.
\]
The latter follows by removing one maximum and one minimum from
the sum of the coordinates and bounding each remaining coordinate
by $L$. The bounds $M-1$ and $(1-M+m)/(n-2)$ coincide at
$M_+:=1+m/(n-1)$.

If $M\ge M_+$, use $b_0\ge M-1$ in
\eqref{eq:r75-null-sum-bound} to obtain
$\sum_i c_i\ge nM-1\ge n-1$. If $M_-\le M\le M_+$,
use $b_0\ge(1-M+m)/(n-2)$. The right side of
\eqref{eq:r75-null-sum-bound} is increasing in $b_0$, and becomes
\[
 F(M)=\frac{1+m}{n-2}+(n-1)M+
              \frac{n-1}{m(n-2)}(M-M^2).
\]
This is a strictly concave quadratic. Its minimum on the stated
interval occurs at an endpoint, where
\[
 F(M_-)=\frac{1+m}{n-1}\left(n+\frac1m\right),
 \qquad F(M_+)=nM_+-1\ge n-1.
\]
The arithmetic--geometric mean inequality gives
\[
 \frac{1+m}{n-1}\left(n+\frac1m\right)
 =\frac{n+1+nm+1/m}{n-1}
 \ge\frac{(\sqrt n+1)^2}{n-1}.
\]
For $n>4$ this last constant is strictly less than $n-1$.
The two cases together prove
$\sum_i c_i\ge(\sqrt n+1)/(\sqrt n-1)$;
\eqref{eq:r75-null-vector-form} proves \eqref{eq:r75-sharp-threshold}.

Suppose equality holds, and write $s=\sqrt n$.
The case $v\ge0$ and the branch $M\ge M_+$ give strictly larger
bounds. Equality in the remaining chain forces
$m=1/s$ and $M=M_-$: strict concavity excludes the interior, and
the other endpoint has value at least $n-1$. Thus
\[
 M=\frac1{s(s-1)},\qquad
 v=(M,\ldots,M,-1/s)
\]
after a coordinate permutation. Indeed the sum of the $n-1$
coordinates other than one minimum is $(n-1)M$, forcing each
to equal its maximum. Now $b_0=M$. Equality in
$v\cdot u\ge-m\sum_i u_i$ places all of $u$ at the unique
negative coordinate, and its value there is $M/m$.
Hence $c=(M,\ldots,M,M+M/m)$. Substituting in
\eqref{eq:r75-null-vector-form} gives precisely
\eqref{eq:r75-extremal-coefficient}.

Conversely the rows $\ell_i$ in that formula are probability
vectors and linearly independent. If $p=\mathbf1^{\mathsf T}/n$
is regarded as a row, their centered rows satisfy
\[
 (\ell_i-p)\cdot(\ell_j-p)
    =\left(1-\frac1{\sqrt n}\right)^2\delta_{ij}.
\]
They are therefore an orthogonal basis of the zero-sum subspace.
Consequently $A_*$ is nonnegative, has total sum one and rank
$n-1$, and
\[
 PA_*P=\frac{(1-1/\sqrt n)^2}{n-1}P
       =\frac{\tau_n}{n}P.
\]
This proves sufficiency and completes the equality characterization.
\end{proof}

\begin{theorem}[The exact two-level nonabelian memory law]
\label{thm:r75-symmetric-return}
For $n\ge5$ and $t>0$, the response \eqref{eq:r75-symmetric-return}
has observable group $S_n$, residual cardinality $n!$, and observable
dimension $1+(n-1)^2$. Its permutation-module bound is
\[
 \nu_{S_n}(V_{b_t})=n(n-1).
\]
At the sharp threshold $\tau_n=(\sqrt n-1)/(\sqrt n+1)$, the
unrestricted positive core and pure-start memory satisfy
\[
 \begin{array}{c|c|c}
 \text{contrast}&\kappa_{S_n}(b_t)&m(b_t)\\\hline
 0<t\le\tau_n&n(n-1)&n(n-1)+1\\
 \tau_n<t\le1&n^2&n^2+1.
 \end{array}
\]
The smaller memory is attained at the threshold itself. At $t=0$
both minima are one.

Every minimum permutation core at $t=\tau_n$ is the disjoint union
of $n-1$ natural orbits. Its probability preparation and binary
decoder, after relabeling the orbit copies and simultaneously
relabeling their natural coordinates, have exactly the form
\begin{equation}\label{eq:r75-threshold-factorizations}
 \lambda_i=\ell_i,\qquad
 D_i=\frac{\ell_i}{(n-1)w_i},\qquad
 \sum_{i=1}^{n-1}w_i=1,\qquad
 w_i\ge\frac{1-1/\sqrt n}{n-1},
\end{equation}
where $\ell_i$ is defined in \eqref{eq:r75-extremal-coefficient},
$w_i$ is the preparation mass on orbit $i$, and $D_i$ is its
first-output decoder written as a row. The second-output decoder
is $1-D_i$.
Conversely every displayed weight vector gives an optimal
factorization. This is a classification of minimum permutation-core
factorizations under the inherited normal form, not a classification
of all stochastic transition realizations before that reduction.
\end{theorem}
\begin{proof}
The real natural permutation module is the direct sum of the
invariant line and the $(n-1)$-dimensional standard representation
$\mathrm{Std}_n$. Its complexification is irreducible on that
zero-sum subspace. One direct proof uses the transpositions: from
a nonconstant vector in an invariant subspace, a vector
$v-(ij)v$ produces a nonzero multiple of $e_i-e_j$, whose orbit
spans the whole zero-sum space. The standard representation is
faithful, because a permutation fixing every difference $e_i-e_j$
fixes all points. Matrix-coefficient orthogonality gives the
nontrivial block
\[
 \widehat\beta_t(\mathrm{Std}_n)
   =\frac{t}{n(n-1)}I_{n-1},
\]
and all other nontrivial irreducible blocks vanish. Therefore the
observable module is
$\boldsymbol1\oplus\mathrm{Std}_n^{\oplus(n-1)}$, of dimension
$1+(n-1)^2$. Equality of two residuals forces equality of their
standard representation matrices because this Fourier block is
invertible; faithfulness then makes the group elements equal.
This proves the claims about the group and residual count.

We next record the subgroup index estimate needed for an unrestricted
memory converse. For $K\le S_n$, let $r$ be the number of its orbits
on $[n]$. The multiplicity of $\mathrm{Std}_n$ in
$\mathbb C[K\backslash S_n]$ is
\[
 m(K)=\dim\mathrm{Std}_n^K=r-1.
\]
If $m(K)>0$, then
\begin{equation}\label{eq:r75-index-estimate}
 [S_n:K]\ge n\,m(K).
\end{equation}
More sharply, if $m(K)>0$ and $K$ is not a full point stabilizer,
\begin{equation}\label{eq:r75-index-extra}
 [S_n:K]\ge n\{m(K)+1\}.
\end{equation}
To prove these, let the natural orbit sizes of $K$ be
$a_1,\ldots,a_r$. Since $K$ is contained in the product of the
symmetric groups on those orbits,
\[
 [S_n:K]\ge\frac{n!}{a_1!\cdots a_r!}.
\]
For $r\ge3$, this is at least $n!/(n-r+1)!$: transferring one
unit from the smaller of two nonunit parts to the larger increases
their factorial product until only one nonunit part remains. The
last lower bound is at least $nr$ for $n\ge5$. Indeed at $r=3$
it is $n(n-1)\ge3n$, and increasing $r$ multiplies the bound by
at least two while $r$ increases by one. If $r=2$ and one part
is a singleton, $K$ lies in a point stabilizer; it is that full
stabilizer of index $n$, or has index at least $2n$. If both
parts have size at least two, the index is at least
$\binom n2\ge2n$. These cases prove both inequalities.
The argument uses the natural $K$-orbits and requires no classification
of low-index subgroups. In particular, an exceptional subgroup of
$S_6$ that is transitive in the natural action has $m(K)=0$ and
does not supply a standard carrier.

A $G$-set containing $V_{b_t}$ must have total standard multiplicity
at least $n-1$. Summing \eqref{eq:r75-index-estimate} over its
standard-carrying orbits gives at least $n(n-1)$ points.
Exactly $n-1$ natural orbits supply this multiplicity, so the module
minimum is $n(n-1)$.

For sharp positive attainment, use the preparations $\ell_i$ of
\eqref{eq:r75-extremal-coefficient} on $n-1$ natural orbits,
each of mass $1/(n-1)$, and use $D_i=\ell_i$ as first-output
decoders. The coefficient matrix is $A_*$. For any permutation
matrix $P_g$, the sum of its entrywise products with a matrix
$A$ will be denoted by $\langle A,P_g\rangle$; in particular
$\lambda P_gD^{\mathsf T}=\langle\lambda^{\mathsf T}D,P_g\rangle$
when the scalar decoder is written as a row.
The identities $PA_*P=(\tau_n/n)P$ and $\sum_{ij}(A_*)_{ij}=1$
imply
\[
 \langle A_*,P_g\rangle
 =\frac1n+\frac{\tau_n}{n}\{\chi_n(g)-1\}
 =\beta_{\tau_n}(g).
\]
Indeed the centered part contributes
$(\tau_n/n)(\chi_n(g)-1)$; the row-constant and column-constant
parts contribute their common group mean $1/n$.
Mixing each scalar decoder with $1/n$ by the factor $t/\tau_n$
realizes every $0\le t\le\tau_n$ on these same points.
The module bound proves optimality for positive $t$ in this interval.

A second, rational construction on the subinterval
$0\le t\le1/(n-1)$ is useful independently. Use $n-1$ natural orbits indexed by
$i=1,\ldots,n-1$, each with initial mass $1/(n-1)$ concentrated
at its point $i$. Give orbit $i$ the scalar decoder
\begin{equation}\label{eq:r75-symmetric-decoder}
 D_i(j)=\frac1n+\frac{(n-1)t}{n}
                 \{\mathbf1_{\{j=i\}}-\mathbf1_{\{j=n\}}\}.
\end{equation}
For $0\le t\le1/(n-1)$ these numbers lie in $[0,1]$.
Since a permutation has exactly one preimage of $n$,
\[
 \frac1{n-1}\sum_{i=1}^{n-1}D_i(g(i))
 =\frac1n+\frac tn
   \sum_{i=1}^{n-1}
    \{\mathbf1_{\{g(i)=i\}}-\mathbf1_{\{g(i)=n\}}\}
 =\frac1n+\frac tn\{\chi_n(g)-1\}=\beta_t(g).
\]
This retains a realization using only rational probabilities whenever
$t$ is rational. Its endpoint is not the sharp threshold when $n\ge5$.

At full contrast, $n$ natural orbits realize $\beta_1$: give
each mass $1/n$, start orbit $i$ at $i$, and use the indicator
decoder of $i$. Hence $\kappa_{S_n}(b_1)\le n^2$.
Suppose, to the contrary, that a factorization exists on fewer
than $n^2$ points. The total standard multiplicity is at least
$n-1$. If any standard-carrying orbit were not natural,
\eqref{eq:r75-index-extra} and the ordinary estimates on the
others would force at least $n(n-1)+n=n^2$ points. Thus all
standard-carrying orbits are natural. There must be exactly
$n-1$ of them: each supplies multiplicity one and costs $n$
points. Every other orbit has total size less than $n$.

Write the factorization as a mixture over its positive-mass
transitive orbits. Since $\beta_1(e)=1$, each component's
decoder equals one on its initial support. On a transitive orbit
of degree $d$ its mean is therefore at least $1/d$.
The natural components have mean at least $1/n$; any additional
orbit here has $d<n$ and would have mean strictly greater than
$1/n$. The mean of $\beta_1$ is exactly $1/n$, so all additional
orbits have zero initial mass and every active natural component
has mean $1/n$. Lemma~\ref{lem:r75-return-rigidity} now makes
each component a deterministic return law. Let $w_i$ be the
aggregate mass of the components starting at natural point $i$.
Evaluating at a transposition $(ij)$ gives
\[
 1-w_i-w_j=\frac{n-2}{n},
 \qquad w_i+w_j=\frac2n.
\]
For $n\ge3$ these equations imply $w_i=1/n$ for every $i$.
Thus at least $n$ positive components are needed, contradicting
the availability of only $n-1$ natural orbits. This proves
$\kappa_{S_n}(b_1)=n^2$.

We next exclude every intermediate positive-core value. Fix $t>0$
and suppose a factorization has fewer than $n^2$ points. The same
index argument gives exactly $n-1$ natural standard-carrying orbits;
the remaining orbits have standard multiplicity zero. Write
$b_t=B_N+B_E$ for their unnormalized response contributions. The
Fourier support of $B_N$ lies in the trivial and standard types,
because these are the only types in a natural permutation module.
The target has the same support. On the other hand, $B_E$ has zero
standard block. Subtracting $B_N$ from $b_t$ shows that all nontrivial
blocks of $B_E$ vanish, so $B_E(g)=v_E$ is a constant nonnegative
vector. Its total mass $w_E=\sum_c v_E(c)$ equals the preparation mass
on the extra orbits. Let $w_N=1-w_E$. Since $t>0$, the nonzero standard
block forces $w_N>0$. Restrict the original preparation and decoder to
the natural orbits, calling them $\lambda_N$ and $D_N$. On those same
$n(n-1)$ points use
\[
 \lambda'_N=\lambda_N/w_N,\qquad
 D'_N=w_ND_N+\mathbf1 v_E.
\]
Every row of $D'_N$ is nonnegative and sums to $w_N+w_E=1$, and
\[
 \lambda'_N P_gD'_N=B_N(g)+v_E=b_t(g).
\]
Thus any positive core below $n^2$ compresses to the module minimum
$n(n-1)$. No intermediate value is possible.

The full-contrast construction realizes every $b_t$ by mixing its
decoder with the constant $1/n$, so $\kappa\le n^2$ everywhere.
It remains to exclude the smaller core for every $t>\tau_n$.
By the absorption argument, any core below $n^2$ would give a
factorization on exactly $n-1$ natural orbits. Write its orbit
weights, probability rows, and scalar decoder rows as
$w_i$, $\lambda_i$, and $D_i$. The matrix
\[
 A=\sum_{i=1}^{n-1}w_i\lambda_i^{\mathsf T}D_i
\]
is entrywise nonnegative and has rank at most $n-1$.
Its response is $\langle A,P_g\rangle=\beta_t(g)$.
The group mean is $(1/n)\sum_{ij}A_{ij}=1/n$, so
$\sum_{ij}A_{ij}=1$. Moreover $PAP=(t/n)P$.
Here is an elementary way to see the last implication without a
choice of Fourier bases. Put $B=A-(t/n)I$. All permutation sums
$\langle B,P_g\rangle$ are equal. Comparing two permutations
that differ only by interchanging the assigned columns of two
distinct rows gives
\[
 B_{ij}+B_{k\ell}=B_{i\ell}+B_{kj}
 \quad(i\ne k,\ j\ne\ell).
\]
Thus $B_{ij}=B_{i1}+B_{1j}-B_{11}$, so $PBP=0$.
Lemma~\ref{lem:r75-singular-coefficient} now forces $t\le\tau_n$.
Together with the matching construction and the absence of
intermediate values, this proves the exact two-level law.
Theorem~\ref{thm:r73-memory-formula}, namely
$m=\min\{n!,1+\kappa\}$, gives the displayed pure-start values since
$n^2+1<n!$ for $n\ge5$. The explicit quadratic threshold is also
consistent with the retained algebraic-threshold theorem; no
quantifier-elimination calculation is needed here.

Finally consider a minimum core at $t=\tau_n$. The index bounds
force exactly $n-1$ natural orbits and no others, and the equality
part of Lemma~\ref{lem:r75-singular-coefficient} forces $A=A_*$
after simultaneous coordinate relabeling. Set
$a=1-1/\sqrt n$ and $b=1/\sqrt n$, so $a>b>0$ for $n\ge5$.
The leading $(n-1)\times(n-1)$ block of $A_*$ is the positive
diagonal matrix $a^2I/(n-1)$. Each nonnegative rank-one summand
$w_i\lambda_i^{\mathsf T}D_i$ can cover at most one of these
diagonal positions: covering two would give a positive off-diagonal
entry. All $n-1$ positions need coverage and there are exactly
$n-1$ summands. Therefore each is covered by a different
positive-weight summand. Relabel that summand by its position $i$.
The zero entries in row and column $i$ then force both
$\lambda_i$ and $D_i$ to be supported on $\{i,n\}$.
Only this summand contributes to the entries $(i,i)$, $(i,n)$,
and $(n,i)$. Their ratios in $A_*$ are
\[
 \frac{(A_*)_{in}}{(A_*)_{ii}}
 =\frac{(A_*)_{ni}}{(A_*)_{ii}}=\frac ba.
\]
Normalization of $\lambda_i$ and the diagonal value
$(A_*)_{ii}=a^2/(n-1)$ give
$\lambda_i=ae_i+be_n=\ell_i$ and
$D_i=\ell_i/\{(n-1)w_i\}$. The condition $D_i\in[0,1]^n$
is exactly $w_i\ge a/(n-1)$ because $a>b$.
Conversely these weights make every decoder valid and give
$w_i\lambda_i^{\mathsf T}D_i=\ell_i^{\mathsf T}\ell_i/(n-1)$.
This proves precisely the claimed classification.
\end{proof}

\paragraph{Prepared quantum realization.}
The family has a finite-dimensional prepared unitary realization,
so it remains inside the experiment framework. On
$\mathbb C^n\otimes\mathbb C^n$ take
\[
 \rho=\frac1n\sum_{i=1}^n|ii\rangle\langle ii|,
 \qquad U_g=P_g^{\mathsf T}\otimes I,
 \qquad E=\sum_{i=1}^n|ii\rangle\langle ii|.
\]
Here $P_g$ is the row-transition permutation matrix, so
$U_g|x,y\rangle=|xg,y\rangle$, consistently with chronological
command multiplication. Then $\operatorname{tr}(E U_g\rho U_g^*)=\chi_n(g)/n$.
The binary effect $E_t=tE+(1-t)I/n$ produces
\eqref{eq:r75-symmetric-return}. Two command types suffice uniformly in
$n$: an $n$-cycle and a transposition of adjacent points generate $S_n$,
since conjugation by the cycle supplies the adjacent transpositions.
They give one common experiment valid for every word in that fixed-size
command alphabet, with a binary terminal outcome set.
The pure maximally entangled preparation
$n^{-1/2}\sum_i|ii\rangle$ gives the same terminal probabilities
with this effect, if a pure preparation is desired. The statement
is a physical realization of the response, not a claim that this
Hilbert-space dimension is minimal or that it establishes a
quantum-versus-classical separation under a common resource metric.

The gap $n$ between the positive core at full contrast and its
module minimum is unbounded. The relative ratio $n/(n-1)$ tends
to one; no growing relative separation is asserted. The exact
transition tends to one as $n$ grows. Its proof combines a general
orbit-multiplicity converse, a sharp positivity inequality with
equality rigidity, and explicit optimal implementation; none of
these counts requires the completion of the wider Theta pipeline.
